\section{Implementation Environment}

The system was developed and tested using the following technology stack:

\begin{table}[htbp]
\centering
\begin{tabular}{l l}
\toprule
\textbf{Component} & \textbf{Technology} \\
\midrule
Stream Broker       & Apache Kafka 7.5 (via Docker Compose) \\
Coordination        & Apache Zookeeper (via Docker Compose) \\
State Store         & Redis 7 \\
API Server          & FastAPI + Uvicorn (Python 3.11+) \\
Python Worker       & kafka-python-ng, mmh3, psutil, redis-py \\
C++ Worker          & C++20, librdkafka, xxHash, CMake + Ninja \\
Dashboard           & HTML5, CSS3, Chart.js, vanilla JavaScript \\
Containerization    & Docker Compose (multi-service orchestration) \\
Testing             & pytest (28 Python tests), Catch2 (C++ tests) \\
\bottomrule
\end{tabular}
\caption{Technology stack used in the Phase~1 implementation.}
\label{tab:tech_stack}
\end{table}

All infrastructure services (Kafka, Zookeeper, Redis) are orchestrated via a single \texttt{docker-compose.yml} file. The C++ worker can optionally be built and run via Docker using a multi-stage \texttt{Dockerfile.worker}.

\section{Interactive Dashboard}

A fully functional, single-page real-time web dashboard was developed to visualize cluster metrics. The dashboard polls the \texttt{/metrics/cluster} endpoint every 2 seconds and renders all data using Chart.js for interactive charting.

\begin{figure}[htbp]
    \centering
    % PLACEHOLDER: Replace with your dashboard screenshot in Overleaf
    % \includegraphics[width=0.9\textwidth]{Figures/dashboard_screenshot.png}
    \framebox[\textwidth]{\rule{0pt}{7cm}\parbox{0.8\textwidth}{\centering\Large\textbf{[PLACEHOLDER]}\\[1em]\normalsize Upload your dashboard screenshot to the \texttt{Figures/} folder in Overleaf and uncomment the \texttt{includegraphics} line above.}}
    \caption{Web dashboard showing live sketch estimates, heavy hitters list, P99 latency charts, and system resource usage.}
    \label{fig:dashboard_screenshot}
\end{figure}

The dashboard interface is organized into the following sections:

\subsection{Live Sketch Estimate Cards}
Three prominent cards display the live answers from each probabilistic sketch:
\begin{itemize}
    \item \textbf{Unique Items} (HyperLogLog): Estimated number of distinct items with relative error percentage.
    \item \textbf{Top Heavy Hitter} (Misra--Gries): The most frequent item in the current sliding window and its share.
    \item \textbf{Top Item Frequency} (Count-Min Sketch): Estimated frequency count with error bound.
\end{itemize}

\subsection{Key Performance Indicator Strip}
Six compact metric boxes display HLL cardinality, CMS total events processed, P99 latency, average CPU usage, average memory usage, and real-time throughput (events per second).

\subsection{Charts}
\begin{itemize}
    \item \textbf{P99 Latency History:} A multi-line time-series chart showing the latency trend per worker over time. Each worker is rendered in a different color with a pulsing ``LIVE'' indicator.
    \item \textbf{CPU and Memory per Worker:} A grouped bar chart showing both CPU and memory usage side by side for each worker.
\end{itemize}

\subsection{Heavy Hitters Table}
A ranked table of the top-12 most frequent items detected by Misra--Gries, with visual percentage share bars for each entry.

\subsection{Query Explorer}
An interactive panel allowing users to test queries directly from the dashboard:
\begin{itemize}
    \item \textbf{Frequency tab:} Enter an item key and retrieve the CMS frequency estimate with error bounds.
    \item \textbf{Cardinality tab:} Retrieve the current HLL distinct count estimate.
    \item \textbf{Heavy Hitters tab:} Retrieve the top-$N$ items from Misra--Gries.
\end{itemize}

\section{Testing and Validation}

\subsection{Unit Test Suite}
A comprehensive test suite of 41 Python unit tests validates the correctness of all Phase~1 and Phase~2 components:

\begin{table}[htbp]
\centering
\begin{tabular}{l c l}
\toprule
\textbf{Component} & \textbf{Tests} & \textbf{Key Assertions} \\
\midrule
Count-Min Sketch     & 6  & Error bounds, no-undercount guarantee, memory \\
HyperLogLog          & 7  & Cardinality accuracy, duplicate handling, precision levels \\
Misra--Gries         & 7  & Heavy hitter detection, sliding-window eviction, sorting \\
Load-Aware Router    & 8  & Scoring formula, node ranking, offline laziness, edge cases \\
Phase 1 Sketch Fixes & 5  & Conservative update, zero-loss migration, decay, HLL/MG resize \\
Phase 2 ML Agents    & 8  & Q-Learning convergence, LinUCB learning, L1 drift, persistence \\
Scatter-Gather \& Monoids & 8 & Monoid laws (assoc/ident), pair election, hedging, unhedged \\
\midrule
\textbf{Total Offline Tests} & \textbf{49} & \textbf{100\% passing offline in $<$0.5 seconds} \\
\bottomrule
\end{tabular}
\caption{Summary of the Python unit test suite (49 tests total, all passing offline without Kafka/Redis).}
\label{tab:test_suite}
\end{table}

All 49 tests pass completely offline without requiring any external infrastructure (no Kafka or Redis needed). The C++ sketches are independently validated via Catch2 unit tests.

\subsection{Sample Query Validation}
The following sample query was executed against a running system to verify end-to-end live RPC frequency estimation:

\begin{lstlisting}[language=bash, caption=Sample live frequency query and response via Redis RPC.]
$ curl "http://localhost:8000/query/frequency?key=item_1"

{
  "key": "item_1",
  "estimated_frequency": 38412,
  "routed_to": "worker-0",
  "router_score": 0.1727,
  "cms_error_bound": 323.31,
  "cms_total_count": 243590,
  "empirical_rel_error": 0.0041,
  "rpc_live": true,
  "worker_metrics": {
    "template": "medium",
    "p95_latency_ms": 0.42,
    "p99_latency_ms": 0.68,
    "memory_percent": 61.4
  }
}
\end{lstlisting}

With a Zipfian distribution ($s = 1.2$), \texttt{item\_1} consistently exhibits the highest estimated frequency, confirming the correctness of the Conservative-Update Count-Min Sketch and live RPC path. Similarly, HLL cardinality estimates for a 10,000-item vocabulary consistently fall within the expected $\pm$1.6\% error range for $p = 12$.

\section{Phase 2: ML-Driven Adaptive Control Evaluation}

In Phase~2, the reactive rule-based controller from Phase~1 was replaced with two online Machine Learning controllers---an \textbf{Online Tabular Q-Learning (RL)} agent and a continuous-feature \textbf{LinUCB Contextual Bandit} baseline---paired with zero-data-loss state-preserving sketch migration, Conservative-Update CMS, real empirical error tracking (WAPE), and L1 Total Variation concept-drift detection.

\subsection{RL and Contextual Bandit Formulations}

\begin{itemize}
    \item \textbf{Online Tabular Q-Learning (Sequential MDP):} Discretises each worker's telemetry into $4 \times 3 \times 3 \times 3 = 108$ states spanning latency (including a proactive ``approaching-SLA'' zone at $70\text{--}100\%$ of SLA), memory pressure, empirical CMS relative error, and active precision template. Actions directly select the target template $a \in \{\text{low}, \text{medium}, \text{high}\}$, allowing single-step jumps during sudden traffic spikes. The agent updates $Q(s, a)$ online via TD(0) Bellman updates with $\epsilon$-greedy decay and a shaped positive-baseline composite reward:
    \begin{equation}
        R_t = 1.0 - \left(w_e \cdot C_{\text{err}} + w_l \cdot C_{\text{lat}} + w_m \cdot C_{\text{mem}} + w_s \cdot \mathbb{I}[\text{switched}]\right) + B_{\text{prec}}
    \end{equation}
    where nonlinear hinge penalties strictly discourage latency ($>200\text{ ms}$) and memory ($>80\%$) SLA violations while $w_s$ suppresses decision flapping.

    \item \textbf{LinUCB Disjoint Contextual Bandit (Continuous-Feature Baseline):} Operates directly on continuous context vectors $x_t \in \mathbb{R}^6$ (bias, normalised latency, memory, empirical error, current template, and L1 drift score) without state discretisation, selecting arms via $p_{t,a} = \hat{\theta}_a^\top x_t + \alpha_{\text{ucb}}\sqrt{x_t^\top A_a^{-1} x_t}$ with online rank-1 ridge updates.
\end{itemize}

\subsection{Five-Configuration Multi-Regime Benchmark}
Using \texttt{scripts/run\_simulation.py}, five configurations were benchmarked across 40,000 Zipfian ($s=1.2$) events on three heterogeneous worker nodes cycling through normal load, accuracy-critical skew with concept drift, bursty latency spikes, and memory pressure.

\begin{table}[H]
\centering
\begin{tabular}{c l r r r r r r}
\toprule
\textbf{\#} & \textbf{Configuration} & \textbf{Reward} & \textbf{SLA Viol\%} & \textbf{P99 (ms)} & \textbf{CMS Err\%} & \textbf{Mem (KB)} & \textbf{Switches} \\
\midrule
1 & Static + Round-Robin (Baseline)     & $+0.5324$ & $41.67\%$ & $111.48$ & $0.962\%$ & $103.5$ & $0$ \\
2 & Rule-Based + Round-Robin (Phase 1)  & $+0.5179$ & $8.75\%$  & $71.37$  & $2.291\%$ & $51.8$  & $21$ \\
3 & Rule-Based + Smart Router (Phase 1) & $+0.5179$ & $8.75\%$  & $56.46$  & $2.465\%$ & $51.8$  & $21$ \\
4 & LinUCB Bandit + Smart Router (Ph.2) & $+0.5466$ & $28.75\%$ & $70.77$  & $1.510\%$ & $90.7$  & $54$ \\
5 & \textbf{Q-Learning RL + Smart (Ph.2)} & \textbf{$+0.5644$} & \textbf{$11.67\%$} & \textbf{$63.39$} & \textbf{$1.913\%$} & \textbf{$75.0$} & \textbf{$16$} \\
\bottomrule
\end{tabular}
\caption{Phase~1 vs. Phase~2 comparative evaluation across 40,000 Zipfian events and 3 heterogeneous nodes under multi-regime traffic bursts, memory pressure, and concept drift.}
\label{tab:phase2_benchmark}
\end{table}

\begin{figure}[H]
    \centering
    \includegraphics[width=0.88\textwidth]{Phase2/fig1_rl_convergence.png}
    \caption{Online policy convergence over 120 multi-regime episodes. Sequential Q-Learning ($\gamma=0.85$) converges to the highest mean reward, outperforming myopic LinUCB, the static baseline, and the Phase~1 rule-based controller.}
    \label{fig:rl_convergence}
\end{figure}

\begin{figure}[H]
    \centering
    \includegraphics[width=0.95\textwidth]{Phase2/fig2_controller_comparison.png}
    \caption{Four-metric comparison across all five configurations: composite Pareto reward, SLA violation rate, empirical CMS WAPE error, and template switching (flapping) count.}
    \label{fig:controller_comparison}
\end{figure}

\begin{figure}[H]
    \centering
    \includegraphics[width=0.86\textwidth]{Phase2/fig3_cu_cms_comparison.png}
    \caption{Standard Count-Min Sketch (Phase~1) versus Conservative-Update Count-Min Sketch (Phase~2) on identical $512 \times 3$ memory under Zipf ($s=1.2$) traffic. Conservative updating reduces WAPE overestimation by more than $15\times$ (from $\approx 6.5\%$ to $\approx 0.4\%$).}
    \label{fig:cu_cms_comparison}
\end{figure}

Key empirical takeaways from Figures~\ref{fig:rl_convergence}--\ref{fig:cu_cms_comparison}:
\begin{enumerate}
    \item \textbf{Superior Pareto Optimality:} Q-Learning achieves the highest composite reward ($+0.5644$), outperforming Static ($+0.5324$), Phase~1 Rule-Based ($+0.5179$), and LinUCB ($+0.5466$).
    \item \textbf{Why Sequential RL Beats Myopic Rules and Bandits:} Fixed rules collapse onto \texttt{low} precision after bursts (yielding $2.465\%$ CMS error and lower utility than static), whereas myopic LinUCB reacts aggressively to continuous feature changes (54 switches, $28.75\%$ SLA violations). Q-Learning's discounted Bellman lookahead ($\gamma=0.85$) and switching penalty deliver \textbf{24\% fewer switches} than Rule-Based (\textbf{70\% fewer} than LinUCB), \textbf{22\% lower CMS error} than Rule-Based, and \textbf{72\% fewer SLA violations} than Static.
    \item \textbf{Over $15\times$ Sketch Accuracy Gain:} Upgrading CMS to Conservative Updating cuts Zipfian overestimation from $\sim 6.5\%$ down to $\sim 0.4\%$ with zero additional memory overhead.
\end{enumerate}

\subsection{Distributed Hedged Scatter-Gather Evaluation under Stragglers}
\label{sec:hedged_evaluation}

To evaluate the horizontal scaling and tail-latency robustness of the Partitioned Hedged Monoid Merge architecture (Section~\ref{sec:hedged_monoid_arch}), we simulated a 4-partition cluster (8 worker nodes total: 4 primary nodes and 4 hot replica nodes) across 1,500 distributed scatter-gather queries using \texttt{scripts/run\_scatter\_gather\_benchmark.py}.

\paragraph{Experimental Setup:}
Each partition's primary node normally serves sub-queries with mean latency $\mu = 2.5\text{ ms}$ ($\sigma = 0.7\text{ ms}$). To emulate real-world production hazards (such as JVM/Python garbage collection pauses, noisy neighbor virtualization stalls, or disk flush stalls), an artificial straggler delay drawn uniformly from $[35\text{ ms}, 75\text{ ms}]$ is injected into primary worker nodes with probability $p_{\text{strag}} = 5.0\%$. We evaluate three query orchestrations:
\begin{enumerate}
    \item \textbf{Unhedged Scatter-Gather:} Standard scatter-gather awaiting all 4 primary partitions.
    \item \textbf{Hedged Scatter-Gather ($T_{\text{hedge}} = 6.0\text{ ms}$):} Duplicate request to replica dispatched if primary exceeds $6.0\text{ ms}$ ($P99$ normal latency).
    \item \textbf{Hedged Scatter-Gather ($T_{\text{hedge}} = 4.5\text{ ms}$):} Duplicate request to replica dispatched if primary exceeds $4.5\text{ ms}$ ($P95$ normal latency).
\end{enumerate}

\begin{table}[htbp]
\centering
\begin{tabular}{l r r r r r r r}
\toprule
\textbf{Configuration} & \textbf{P50 (ms)} & \textbf{P90 (ms)} & \textbf{P95 (ms)} & \textbf{P99 (ms)} & \textbf{Max (ms)} & \textbf{Overhead} \\
\midrule
Unhedged Scatter-Gather       & $3.37$ & $53.89$ & $63.89$ & $73.00$ & $74.90$ & $0.00\%$ \\
Hedged ($T_{\text{hedge}} = 6.0$ ms) & $3.37$ & $8.53$  & $9.14$  & $9.99$  & $11.08$ & $4.85\%$ \\
\textbf{Hedged ($T_{\text{hedge}} = 4.5$ ms)} & \textbf{$3.37$} & \textbf{$7.03$}  & \textbf{$7.64$}  & \textbf{$8.49$}  & \textbf{$9.58$}  & \textbf{$5.65\%$} \\
\bottomrule
\end{tabular}
\caption{Distributed scatter-gather tail latency comparison across 1,500 queries over 4 partitions with 5\% injected worker straggler stalls.}
\label{tab:hedged_benchmark}
\end{table}

\begin{figure}[H]
    \centering
    \includegraphics[width=0.95\textwidth]{Phase2/fig4_hedged_scatter_gather.png}
    \caption{Distributed Scatter-Gather Tail Latency under Stragglers. (a) Cumulative Distribution Function (CDF) showing that unhedged queries suffer a long straggler tail stretching past $74\text{ ms}$, whereas hedged requests sharply truncate the distribution at $\approx 8.5\text{ ms}$. (b) Tail latency percentiles showing an $>8.5\times$ reduction in P99 latency with only 5.65\% duplicate request overhead.}
    \label{fig:hedged_scatter_gather}
\end{figure}

\paragraph{Discussion and Key Takeaways:}
\begin{itemize}
    \item \textbf{Catastrophic Amplification in Unhedged Scatter-Gather:} Even though each individual worker experiences a straggler stall only 5\% of the time, the probability that at least one of the 4 partitions encounters a straggler is $1 - (1 - 0.05)^4 = 18.55\%$. Consequently, in the unhedged configuration, nearly one in every five queries stalls, dragging P90 out to $53.89\text{ ms}$ and P99 out to $73.00\text{ ms}$.
    \item \textbf{Drastic Tail Latency Truncation ($>8.5\times$ Speedup):} With Hedged Scatter-Gather set at $T_{\text{hedge}} = 4.5\text{ ms}$ (the empirical P95 threshold), P90 latency drops from $53.89\text{ ms}$ to $7.03\text{ ms}$ ($7.7\times$ faster), and P99 latency drops from $73.00\text{ ms}$ to $8.49\text{ ms}$ (\textbf{$8.6\times$ faster}). Maximum observed query latency across all 1,500 queries drops from $74.90\text{ ms}$ down to $9.58\text{ ms}$.
    \item \textbf{Negligible Network Overhead:} Because hedged duplicate requests are only triggered when the primary worker fails to reply within the normal operational window, only 339 duplicate sub-queries were dispatched out of 6,000 total sub-queries, representing a modest **5.65\% network overhead**.
    \item \textbf{Sub-Millisecond Monoid Merge:} Merging the partial Count-Min matrices, HyperLogLog registers, and Misra--Gries heavy-hitter summaries across all 4 partitions took an average of $38\ \mu\text{s}$, confirming that algebraic Monoid merges introduce near-zero CPU overhead at query time.
\end{itemize}
